\documentclass[quick,con]{debate}
\motion{对知识网红的崇拜让我们离真知更远（反方）}
\begin{document}
\debatetitle
\begin{thesisbox}[一句话立论]崇拜改变的不是我们接触多少知识，而是我们凭什么相信它——当信念的支点从证据换成一个人，我们拿到的只是结论，丢掉的是理解和纠错，所以离真知更远。\end{thesisbox}
\section{定义与判准（原话）}
\begin{fields}
\field{知识网红}{在网上传播知识、并因此获得大量关注的创作者（与对方同一版，不玩文字游戏）。
}
\field{崇拜}{明显强于一般关注的推崇与信任，表现为持续追随、\textbf{优先采信其观点}。
}
\field{真知}{正确而有理据的认识，三层——内容对，知道自己凭什么认为它对，错了能改。
}
\field{判准}{跟“同样在看这些人、但没有崇拜”的自己比，崇拜这一层在接触、正确性、理据、纠错四项上的净额是往前还是往后。\textbf{我方要证明}理据与纠错上的损耗大于接触增量；\textbf{对方要证明}增量大于损耗。
}
\field{两条对称禁令}{短视频媒介的毛病不能记在崇拜头上；知识网红的功劳同样不能记在崇拜头上。
}
\field{对方提偏向定义或判准时的 30 秒口径}{①“崇拜”若只是“很喜欢”，这道题不需要辩；词典给的是“尊敬钦佩”，重量在于把对方放在自己之上。②我方承认崇拜带来了接触，“更”字要求算净额，不是有没有好处。③主观起点前移，实际起点没动，而人是按主观起点行动的。
}
\end{fields}
\section{我方论点}
\begin{longtable}{@{}L{\dimexpr0.050\linewidth-1.50\tabcolsep}L{\dimexpr0.099\linewidth-1.50\tabcolsep}L{\dimexpr0.433\linewidth-1.50\tabcolsep}L{\dimexpr0.418\linewidth-1.50\tabcolsep}@{}}
\toprule \thead{标签} & \thead{一句话} & \thead{核心数据 / 学理 / 案例} & \thead{出处}\\\midrule\endhead
\textbf{用人替证据} & 崇拜改变的不是接触多少，而是凭什么相信 & 学理：Sperber 等 2010 认知警惕分信源层与内容层，崇拜把内容层外包出去；Hardwig 1985 的条件是“有好的理由相信对方有好的理由”，崇拜把它换成了讲得生动。数据：166 条疫情短视频 22\% 中度、7\% 高度错误信息——对错混在一起，“信谁”替代不了“信什么”。案例：AI 合成张文宏卖蛋白棒 1266 件 & \emph{Mind \& Language} 25(4)；\emph{J. Philosophy} 82(7)；Baghdadi et al. 2023, \emph{JAMIA Open}；央视网 2024-12-18\\
\textbf{关掉除错开关} & 错觉是人性，崇拜是那只按住除错开关的手 & 学理：Rozenblit \& Keil 2002 解释深度错觉——被要求逐步解释，自评当场下跌。数据：Fisher/Goddu/Keil 2015 的 9 项实验，搜过就自评更懂，哪怕没搜到；Salomon 1984 对 124 名六年级学生，觉得容易就投入更少、成绩更低。案例：2023 年 6 月“报新闻学我把他打晕”，出来讲理由的是重庆大学教授 & \emph{Cognitive Science} 26(5)；\emph{JEP: General} 144(3)；ERIC EJ306057；观察者网 2023-06-18\\
\bottomrule\end{longtable}
\section{对方论点 → 一句话反驳}
\begin{longtable}{@{}L{\dimexpr0.230\linewidth-1.33\tabcolsep}L{\dimexpr0.365\linewidth-1.33\tabcolsep}L{\dimexpr0.404\linewidth-1.33\tabcolsep}@{}}
\toprule \thead{对方会说} & \thead{我方一句话回} & \thead{对方追问时}\\\midrule\endhead
门票效应：崇拜把一次性接触变成持续关注 & 崇拜在哪一步让人更懂得他相信的东西为什么是对的？ & 要转化率：从入口到深入的比例是多少？没有数字就是一句假设\\
科学素质 14.14\%、互联网首选 58.3\% & 这条数据里没有“崇拜”这个变量 & 对方自己承认它是相关不是因果，那今天还剩什么是关于崇拜的\\
崇拜给错误一个可以被撤回的地址 & 那个地址上的人自己都说，投诉像蝗灾一样打不完 & 声誉惩罚的是人设崩塌和带货翻车，不指向准确性\\
汪品先院士愿意留在公共场域 & 我方也鼓掌，但他自己说这是“被上网”；院士的可信度来自研究，不来自弹幕 & 那是关注度不是崇拜；请证明“崇拜”那一部分的独立贡献\\
Hardwig：依赖专家是理性的 & 他的条件是“有好的理由相信对方有好的理由”，崇拜把理由换成了讲得好听 & 合理做法是核对资质、看同行共识、看他是否在自己领域内说话——崇拜者最不做这三件事\\
\bottomrule\end{longtable}
\section{必问与必答}
\begin{points}{必问 （对方不答就点名、重复、不换题）}
\item 崇拜在哪一步让人更懂得他相信的东西为什么是对的？（一辩稿抛出的问题）
\item 从入口到深入的比例是多少？
\item 复制一张脸就带走整份信任，这说明信任长在哪里？
\item 声誉惩罚的是讲错知识点，还是人设崩塌？
\item 那几百万转发的人里，有几个去查过数据？
\end{points}
\begin{points}{必答 （15 秒正面回答，不许回避）}
\item “你方数据的自变量是崇拜吗？”→ 不是，我方主动说明。所以今天比机制，我方的机制走完了三步。
\item “起点前移算不算更近？”→ 主观起点前移，实际起点没动，而人是按主观起点行动的。
\item “不依赖专家难道全靠自己验证？”→ 依赖没问题，依赖要有理由；崇拜把理由换成了魅力。
\item “你方是不是在骂知识网红？”→ 不是。去掉“崇拜”两个字，我方一句话都说不出来。
\end{points}
\section{自由辩战场概要}
\begin{longtable}{@{}L{\dimexpr0.212\linewidth-1.50\tabcolsep}L{\dimexpr0.222\linewidth-1.50\tabcolsep}L{\dimexpr0.202\linewidth-1.50\tabcolsep}L{\dimexpr0.364\linewidth-1.50\tabcolsep}@{}}
\toprule \thead{战场} & \thead{开场问题} & \thead{目标} & \thead{收束语}\\\midrule\endhead
一·机制那一步（前 90 秒） & 崇拜在哪一步让人更懂得“为什么”？ & 让评委听见对方始终没有回答 & “正方描述了一扇门，却没告诉我们门后面有什么。”\\
二·转化率（中 90 秒） & 一百个崇拜者里有几个真的深入学习了？ & 让评委听见这个数字一次都没给过 & “他们数的是走到门口的人，我们问的是走进去的人。”\\
三·收束（后 60 秒） & 不开新战场 & 重复一句话立论与必问 & “错觉是人性，崇拜是按住除错开关的手。”\\
\bottomrule\end{longtable}
\section{如果……就……}
\begin{contingency}
\ifthen{对方把崇拜淡化成“很喜欢”，或指控我方打稻草人}{“如果崇拜只是很喜欢，这道题不需要辩”；我方全场只用“优先采信”四个字，这四个字您方也承认。}
\ifthen{对方要我方量化损耗}{“我方给不出比例，您方也给不出转化率；双方都在机制层，那就比机制。”}
\ifthen{用人替证据被打穿，或对方回避必问}{收缩到“关掉除错开关”，用判准第三、四项解释为什么仍占优；对方回避就点名、重复、不换题。}
\end{contingency}
\end{document}
